\documentclass[a4paper,12pt]{article}
\usepackage[spanish]{babel}
\usepackage{amsmath}
\usepackage{graphicx}
\usepackage{multicol}
\usepackage{geometry}
\usepackage{tikz}
\usepackage{fancyhdr}

\geometry{top=1cm, bottom=1cm, left=1cm, right=1cm}
\pagestyle{empty}

\begin{document}

% ==== Evaluación 1 ====

\begin{center}
    \Large \textbf{Evaluación de Termodinámica y Máquinas Térmicas A} \\
    \normalsize Duración: 80 minutos \\
    Alumno/a: \underline{\hspace{6cm}} \hspace{0.5cm} Fecha: \underline{\hspace{2cm}}
\end{center}

\vspace{0.3cm}

\begin{multicols}{2}
\section*{Problemas}
\begin{enumerate}

\item \textbf{Trabajo de compresión en contexto real} \\
En una planta de producción de alimentos, un compresor toma $0{,}5\,\text{m}^3$ de aire a presión atmosférica y lo comprime hasta $0{,}02\,\text{m}^3$ a una presión final de $3\,\text{atm}$. Durante este proceso se realiza un trabajo de compresión de $8000\,\text{J}$. Calcule el trabajo de circulación total.

\vspace{0.3cm}

\item \textbf{Primer principio} \\
Un gas realiza un trabajo de expansión de $2562\,\text{kgm}$, manteniéndose constante su energía interna. Calcule el calor entregado al sistema, usando únicamente el primer principio de la termodinámica.

\vspace{0.3cm}

\item \textbf{Trabajo sobre gas } \\
Calcule el trabajo necesario para introducir $1{,}5\,\text{m}^3$ de gas en un reactor que opera a una presión constante de $6\,\text{atm}$. Explique brevemente la fórmula utilizada.

\vspace{0.3cm}

\item \textbf{Variación de energía interna} \\
Un gas es comprimido mediante un émbolo de área transversal $A = 0{,}01\,\text{m}^2$ que se desplaza $d = 0{,}3\,\text{m}$ bajo presión constante de $2\,\text{atm}$. Durante el proceso se transfieren $5\,\text{kcal}$ al ambiente. Calcule la variación de energía interna del gas.

\vspace{0.3cm}

\item \textbf{Trabajo en almacenamiento de gas} \\
En una instalación de almacenamiento, se desea agregar $1\,\text{m}^3$ de gas a un tanque que ya se encuentra a una presión de $5\,\text{atm}$. Determine el trabajo necesario para efectuar esta tarea.

\end{enumerate}
\end{multicols}

\vspace{1cm}
\hrule
\vspace{1cm}

% ==== Evaluación 2 ====

\begin{center}
    \Large \textbf{Evaluación de Termodinámica y Máquinas Térmicas B} \\
    \normalsize Duración: 80 minutos \\
    Alumno/a: \underline{\hspace{6cm}} \hspace{0.5cm} Fecha: \underline{\hspace{2cm}}
\end{center}

\vspace{0.3cm}

\begin{multicols}{2}
\section*{Problemas}
\begin{enumerate}

\item \textbf{Trabajo externo en sistema cerrado} \\
Un cilindro con un pistón contiene un gas a una presión de $3\,\text{atm}$. Se desea introducir $0{,}5\,\text{m}^3$ de gas en su interior. ¿Cuánto trabajo se necesita realizar para lograrlo?

\vspace{0.3cm}

\item \textbf{Trabajo de circulación} \\
Un compresor recibe $1\,\text{m}^3$ de aire a presión atmosférica y lo entrega comprimido a $0{,}025\,\text{m}^3$ y $3\,\text{atm}$, realizando un trabajo de compresión de $13500\,\text{kgm}$. Determine el trabajo de circulación total.



\vspace{0.3cm}

\item \textbf{Trabajo en sistema de refrigeración} \\
Un sistema de refrigeración emplea un compresor para mantener una presión constante de $2\,\text{atm}$. Si se desea introducir $2\,\text{m}^3$ de gas, ¿cuál es el trabajo que debe realizar el compresor?

\vspace{0.3cm}

\item \textbf{Trabajo efectuado por el gas} \\
Un gas contenido en un cilindro tiene un volumen inicial $v_i = 0{,}025\,\text{m}^3$ y presión $p_i = 2{,}5\,\text{atm}$. Al calentarlo, su volumen aumenta a $v_f = 0{,}175\,\text{m}^3$ manteniéndose constante la presión. Calcule el trabajo efectuado por el gas.

\vspace{0.3cm}

\item \textbf{Trabajo sobre un gas en laboratorio} \\
En un laboratorio, un recipiente sellado contiene gas a una presión de $4\,\text{atm}$. Para un experimento se introduce $0{,}8\,\text{m}^3$ de gas adicional. Determine el trabajo requerido.

\end{enumerate}
\end{multicols}

\end{document}
